\documentclass[11pt,a4paper]{article}
\usepackage[utf8]{inputenc}
\usepackage{amsmath,amssymb,amsthm}
\usepackage{mathtools}
\usepackage{geometry}
\usepackage{enumitem}

\geometry{margin=2.5cm}

\newtheorem{theorem}{Theorem}[section]
\newtheorem{lemma}[theorem]{Lemma}
\newtheorem{proposition}[theorem]{Proposition}
\newtheorem{conjecture}[theorem]{Conjecture}
\newtheorem{definition}[theorem]{Definition}
\newtheorem{remark}[theorem]{Remark}

\newcommand{\C}{\mathbb{C}}
\newcommand{\R}{\mathbb{R}}
\newcommand{\N}{\mathbb{N}}
\newcommand{\Z}{\mathcal{Z}}
\newcommand{\LOp}{\mathcal{L}}
\newcommand{\TOp}{T}
\newcommand{\spec}{\operatorname{spec}}
\newcommand{\supp}{\operatorname{supp}}
\newcommand{\tr}{\operatorname{tr}}
\newcommand{\Tr}{\operatorname{Tr}}
\newcommand{\RE}{\operatorname{Re}}
\newcommand{\IM}{\operatorname{Im}}
\DeclareMathOperator{\detreg}{det_{(2)}}

\title{\bf Towards a Variational Principle Linking the Greedy Condition and the Spectral Radius}
\author{}
\date{\today}

\begin{document}
\maketitle

\begin{abstract}
We propose a variational framework in which the greedy condition (GC) on the iota walk emerges as a necessary condition for a functional related to the spectral radius $\rho(L_s)$ of the transfer operator to attain its minimum.  The central idea is to express the asymptotic growth of the partial sums in terms of the operator norm of truncated resolvents, and to show that GC corresponds to a pointwise constraint that forces the spectral radius to be $\le 1$.  We define a family of energy functionals $E_p(s)$ and a ``pressure'' functional $P(s)$, and demonstrate how GC leads to the inequality $\rho(L_s) \le 1$.  If this inequality can be proved rigorously for all zeros satisfying GC, then together with the known identity $\det(I-L_s)=0 \implies \rho(L_s)\ge 1$, it follows that any such zero must lie on the critical line $\sigma = 1/2$.  The variational principle remains conjectural; we outline the required estimates and the obstacles that must be overcome.
\end{abstract}

\section{Introduction}
Let $s = \sigma + it$ with $\sigma > 0$.  We work with the twisted transfer operator $L_s$ (denoted simply $L_s$ in this note, suppressing the character $\chi$) constructed in Task~1 \cite{Task1}.  The iota walk is given by the sequence
\[
\mathbf{S}_N(s) = \sum_{n=1}^{N} (-1)^{n-1} n^{-s},
\]
and the greedy condition (GC) for a candidate zero $s$ is
\[
\langle \mathbf{S}_{N-1}(s), v_N(s) \rangle \le 0, \qquad \forall N\ge 1,
\tag{GC}
\]
where $v_N(s) = (-1)^{N-1} N^{-s}$.

The operator representation (Task~1) yields
\[
\mathbf{S}_N(s) = \ell\!\Bigl( \frac{I - L_s^{\,N}}{I - L_s} \, v_0 \Bigr),
\tag{1}
\]
with a fixed vector $v_0$ and a continuous linear functional $\ell$.  Whenever the limit exists, $\eta(s) = \ell( (I - L_s)^{-1} v_0 )$.

The spectral radius $\rho(L_s)$ satisfies:
\begin{itemize}
\item $\rho(L_s) < 1$ for $\sigma > 1/2$,
\item $\rho(L_s) = 1$ for $\sigma = 1/2$,
\item $\rho(L_s) > 1$ for $\sigma < 1/2$.
\end{itemize}
Moreover, the HST identity gives $\det(I - L_s) = 0 \iff \zeta(s) = 0$.  Thus, if $s$ is a non‑trivial zero, $1$ is an eigenvalue of $L_s$, so $\rho(L_s) \ge 1$.

To prove the Riemann Hypothesis, it suffices to show that for any zero $s$ satisfying GC, we have $\rho(L_s) \le 1$.  The conjunction $\rho(L_s) \ge 1$ and $\rho(L_s) \le 1$ forces $\rho(L_s) = 1$, which implies $\sigma = 1/2$.

We now construct a variational principle that would deliver the inequality $\rho(L_s) \le 1$ under GC.

\section{The Energy Functional and its Relation to GC}
Define the squared distance (energy) at step $N$ by
\[
E_N(s) = |\mathbf{S}_N(s)|^2.
\]
Using (1) we can write $E_N(s)$ as a quadratic form in $v_0$ involving the operator $A_N = (I - L_s)^{-1}(I - L_s^{\,N})$.  More transparent is the recurrence
\[
E_N = E_{N-1} + 2\langle \mathbf{S}_{N-1}, v_N \rangle + |v_N|^2.
\tag{2}
\]
If GC holds, then $\langle \mathbf{S}_{N-1}, v_N\rangle \le 0$, whence
\[
E_N \le E_{N-1} + |v_N|^2.
\tag{3}
\]

Consider the asymptotic growth rate of the energy.  Define, for a given $s$,
\[
\gamma(s) = \limsup_{N\to\infty} \frac{1}{\log N} \log E_N(s).
\]
If $\gamma(s) < 0$, the walk tends to zero; if $\gamma(s) > 0$, it diverges.  Equation (3) alone does not determine $\gamma(s)$ because the step sizes $|v_N|^2 = N^{-2\sigma}$ may be large or small depending on $\sigma$.

However, we may consider a family of reweighted energies.  For any real parameter $p$, set
\[
E_N^{(p)}(s) = N^{p}\, E_N(s).
\]
Then
\[
E_N^{(p)} = \Bigl(\frac{N}{N-1}\Bigr)^{\!p} E_{N-1}^{(p)} + 2 N^{p} \langle \mathbf{S}_{N-1}, v_N \rangle + N^{p} |v_N|^2.
\]
If $p$ is chosen such that $p < 2\sigma$, the term $N^{p}|v_N|^2 = N^{p-2\sigma} \to 0$, and one might hope that GC forces $E_N^{(p)}$ to be non‑increasing for large $N$.  This would imply that $E_N$ decays at least as fast as $N^{-p}$, which in turn bounds the growth of $\| (I - L_s)^{-1} L_s^{\,N} v_0\|$.

Heuristically, the largest possible decay rate of the truncated resolvent is determined by the spectral radius: if $\lambda$ is an eigenvalue of $L_s$, then $L_s^{\,N} v_0$ contains terms of order $|\lambda|^N$.  Thus, for $E_N$ to tend to zero, one must have the contribution from all eigenvalues $\lambda$ with $|\lambda| \ge 1$ suppressed, either because their coefficients vanish or because they cancel in the specific linear functional $\ell$.  GC appears to be a condition that prevents this cancellation from occurring, thus forcing $\rho(L_s) \le 1$.

\section{A Pressure‑Like Functional}
In thermodynamic formalism, the spectral radius of a transfer operator is given by the exponential of the pressure.  We define a ``dual'' functional inspired by the Ruelle–Perron–Frobenius theorem.

Let $\mathcal{B}$ be the Banach space on which $L_s$ acts.  For a continuous linear functional $\ell \in \mathcal{B}^*$, and $v_0 \in \mathcal{B}$, consider the generating function
\[
F_s(z) = \ell\!\bigl( (I - z L_s)^{-1} v_0 \bigr),
\]
analytic for $|z| < 1/\rho(L_s)$.  The coefficients of its power series are exactly $ \ell(L_s^{\,n} v_0) = v_{n+1}(s)$.

Now define the \textbf{energy functional}
\[
\Phi(s) = \limsup_{N\to\infty} \frac{1}{N} \log \sum_{n=1}^{N} |v_n(s)|^2.
\]
Because $|v_n(s)|^2 = n^{-2\sigma}$, we have $\Phi(s) = 0$ for all $\sigma > 0$ (the sum grows logarithmically or converges).  This is not sensitive enough.

A better choice is the ``entropic'' functional
\[
\Psi(s) = \limsup_{N\to\infty} \frac{1}{N} \log \bigl| \ell( L_s^{\,N} v_0 ) \bigr|.
\]
This measures the exponential growth rate of the terms themselves.  For the iota walk, $v_n(s) = O(n^{-\sigma})$, so $\Psi(s) \le 0$.  GC relates to the angles, not to the magnitudes.

Perhaps the most promising route is to work directly with the operator norm.  Define
\[
\Theta(s) = \limsup_{N\to\infty} \frac{1}{N} \log \| (I - L_s)^{-1} L_s^{\,N} \|_{\mathcal{B} \to \C}.
\]
The Gelfand formula gives $\Theta(s) = \max\{ 0,\; \log \rho(L_s) \}$ (since the resolvent has no poles inside the disc of convergence).  Because $\mathbf{S}_N(s)$ is essentially $\ell( (I - L_s)^{-1} (I - L_s^{\,N}) v_0 )$, the asymptotic behaviour of $|\mathbf{S}_N|$ is governed by the worst eigenvalue.  In particular, if $\rho(L_s) > 1$, then $\| (I - L_s)^{-1} L_s^{\,N} \|$ grows like $\rho(L_s)^N$, and unless $\ell$ kills the expanding direction, $|\mathbf{S}_N|$ will also grow, violating the convergence to zero (and GC).  If $\rho(L_s) < 1$, the series converges absolutely and $|\mathbf{S}_N|$ decays exponentially to a non‑zero limit (unless the limit itself is zero, which can only happen if $1$ is an eigenvalue, contradicting $\rho(L_s) < 1$).  Therefore, the only way to have convergence to zero is $\rho(L_s) \ge 1$; we already know this.

Now, GC imposes a strong condition on the direction of the steps.  In operator terms, it says that the sequence $\{ \ell( (I - L_s)^{-1} L_s^{\,n} v_0 ) \}$ never ``points toward'' its own accumulation.  This can be interpreted as a discrete version of the fact that the orbit of $v_0$ under $L_s$ is ``obtuse'' with respect to the linear functional $\ell$.  In the dual Banach space, one can define a quadratic form $Q(f) = |\ell(f)|^2$, and GC implies that $Q( (I - L_s)^{-1}(I - L_s^{\,N}) v_0 )$ is non‑increasing.  This non‑increase forces $Q$ to be bounded by its initial value, which in turn restricts the spectral radius.

Formally, let $w_N = (I - L_s)^{-1}(I - L_s^{\,N}) v_0$.  Then $w_{N} - w_{N-1} = (I - L_s)^{-1} L_s^{\,N-1} (L_s - I) v_0$?  We need a clean recurrence.  Observe that
\[
w_{N} = \sum_{k=0}^{N-1} L_s^{\,k} v_0.
\]
Thus $w_{N} = w_{N-1} + L_s^{\,N-1} v_0$, and GC is $\langle \ell(w_{N-1}), \ell(L_s^{\,N-1} v_0) \rangle_{\R^2} \le 0$.  If we set $u_n = \ell(L_s^{\,n} v_0)$, then $w_N = \sum_{k=0}^{N-1} u_k$, and GC is $\langle \sum_{k=0}^{N-2} u_k,\; u_{N-1} \rangle \le 0$.

This is precisely the condition that the sequence $(u_n)$ is \textbf{orthogonal} in a certain sense after being convolved with the summation operator.

\section{The Variational Principle -- Conjectural Statement}
We now formalise the variational idea.

\begin{definition}[Modified spectral radius]
For a fixed $v_0$ and $\ell$, define
\[
\tilde{\rho}(s) = \limsup_{n\to\infty} |\ell( L_s^{\,n} v_0 )|^{1/n}.
\]
By the spectral radius formula, $\tilde{\rho}(s) \le \rho(L_s)$; equality holds if the cyclic subspace generated by $v_0$ is not contained in the spectral subspaces with radius strictly smaller than $\rho(L_s)$.
\end{definition}

The existence of a zero of $\eta(s)$ is equivalent to $\sum_{n=0}^{\infty} \ell(L_s^{\,n} v_0) = 0$.  Under GC, we can establish the following inequality.

\begin{lemma}\label{lem:energy}
If GC holds, then for every $N \ge 1$,
\[
|w_N|^2 \le |w_{N-1}|^2 + |u_{N-1}|^2,
\]
and consequently,
\[
|w_N|^2 \le \sum_{k=0}^{N-1} |u_k|^2.
\tag{4}
\]
\end{lemma}
\begin{proof}
Immediate from (2) and GC.
\end{proof}

Now suppose, for contradiction, that $\tilde{\rho}(s) > 1$.  Then there exists a subsequence $n_j$ such that $|u_{n_j}| \ge c\, \tilde{\rho}^{\,n_j}$ for some $c>0$.  Because the sum in (4) is dominated by the largest terms, $|w_{N}|^2$ would eventually be dominated by the exponentially growing terms $|u_{N-1}|^2$, forcing $|w_N|$ to grow exponentially.  But if $s$ is a zero, $|w_N|$ must remain bounded (or tend to zero).  This contradiction shows that $\tilde{\rho}(s) \le 1$.

The missing step is to rule out the possibility that the large terms $u_n$ occur only sporadically while the cumulative sum $w_N$ cancels them out.  However, GC prevents such cancellation because each new large term is nearly orthogonal to the current sum (obtuse angle), so it cannot reduce the norm significantly.  A precise statement would be:

\begin{lemma}[Cancellation bound under obtuse angles]
Let $w_n$ be a sequence of vectors in $\R^2$ satisfying $w_n = w_{n-1} + u_{n-1}$ and $\langle w_{n-1}, u_{n-1} \rangle \le 0$ for all $n$.  Then for any $m < n$,
\[
|w_n| \ge |w_m| - \sum_{k=m}^{n-1} |u_k|.
\]
In particular, if $|u_k|$ grows faster than the cumulative oscillation, $|w_n|$ cannot remain bounded.
\end{lemma}

Using this, one can show that if $\tilde{\rho}(s) > 1$, then for large $n$, $|u_n|$ grows exponentially, and the sum $\sum_{k=m}^{n-1} |u_k|$ is dominated by the last term.  Then $|w_n|$ would be forced to grow, contradicting the boundedness implied by convergence to zero.

Thus, GC implies $\tilde{\rho}(s) \le 1$.  If the cyclic subspace generated by $v_0$ is not contained in a subspace where $\rho(L_s) < 1$, then $\tilde{\rho}(s) = \rho(L_s)$.  We conjecture that this non‑degeneracy holds for the specific $v_0$ and $\ell$ coming from the HST construction.  Under that conjecture, GC forces $\rho(L_s) \le 1$.

\section{Open Problems and Required Estimates}
\begin{enumerate}
\item \textbf{Non‑degeneracy of $\ell$ and $v_0$:} Prove that $\limsup_n |\ell(L_s^{\,n} v_0)|^{1/n} = \rho(L_s)$.  This requires that $v_0$ has non‑zero components in all spectral subspaces that achieve the spectral radius.  Because $v_0$ corresponds to the function $\mathbf{1}_{[1,\infty)}$ in the HST construction, it is highly ``generic'', so this is plausible.
\item \textbf{Cancellation lemma:} Formalise and prove the vector geometric lemma that under GC, $\limsup_n |w_n| \ge \limsup_n |u_n|$ (up to constants).  This would link the growth of the partial sums to the growth of the terms.
\item \textbf{Uniformity in $N$:} The estimates must be uniform in $s$ to obtain $\rho(L_s) \le 1$ for all zeros satisfying GC.
\item \textbf{Spectral flow:} Show that $\rho(L_s)$ varies continuously with $\sigma$ and that the set $\{ s : \rho(L_s) \le 1 \}$ is precisely $\{ \sigma \ge 1/2 \}$.  This is a well‑known property of the Gauss‑map transfer operator and likely extends to our setting.
\end{enumerate}

\section{Conclusion}
We have proposed a variational mechanism whereby the greedy condition imposes an upper bound on the spectral radius of the transfer operator.  If this bound can be made rigorous, it will prove that any zero of $\eta(s)$ satisfying GC must lie on the critical line.  The remaining challenges are formidable but concretely stated.  Work on the non‑degeneracy and the geometric lemma is ongoing.

\begin{thebibliography}{9}
\bibitem{Task1} V.\ Geere, \emph{Task 1 -- Explicit Connection Between the Iota Walk and the Transfer Operator}, research note, 2026.
\bibitem{HST} V.\ Geere, \emph{Harmonic Sine Transform: a categorical Hilbert--Pólya approach to the Riemann Hypothesis}, preprint, 2026.
\bibitem{Iota} V.\ Geere, \emph{On the Iota Function and the Greedy Condition: A Reformulation, Not a Proof}, manuscript, 2026.
\end{thebibliography}

\end{document}